% time_machine_nogo.tex
% Time, Computation, and Novelty: A No-Go Theorem from Petz Recovery
% For submission to Physical Review Letters
% Author: Sheng-Kai Huang
% Compile with: pdflatex time_machine_nogo.tex && pdflatex time_machine_nogo.tex && pdflatex time_machine_nogo.tex

\documentclass[prl,twocolumn,superscriptaddress,longbibliography,floatfix]{revtex4-2}

\usepackage{amsmath}
\usepackage{amssymb}
\usepackage{amsfonts}
\usepackage{amsthm}
\usepackage{bm}
\usepackage{hyperref}
\usepackage{booktabs}
\usepackage{graphicx}

\newtheorem{theorem}{Theorem}
\newtheorem{corollary}[theorem]{Corollary}
\newtheorem*{proposition}{Proposition}

% Convenience macros (matching Paper 1)
\newcommand{\kB}{k_{\mathrm{B}}}
\newcommand{\Tr}{\mathrm{Tr}}
\newcommand{\id}{\mathrm{id}}
\newcommand{\Petz}{\mathcal{R}}
\newcommand{\chan}{\mathcal{N}}
\newcommand{\hilb}{\mathcal{H}}
\newcommand{\code}{\mathcal{C}}
\newcommand{\KL}{D}
\newcommand{\ket}[1]{|#1\rangle}
\newcommand{\bra}[1]{\langle#1|}
\newcommand{\braket}[2]{\langle#1|#2\rangle}
\newcommand{\op}[2]{|#1\rangle\langle#2|}
\newcommand{\abs}[1]{\left|#1\right|}
\newcommand{\norm}[1]{\left\|#1\right\|}
\newcommand{\tauasym}{\tau}
\newcommand{\taurot}{\tilde{\tau}}

\begin{document}

\title{Time, Computation, and Novelty: A No-Go Theorem from Petz Recovery}

\author{Sheng-Kai Huang}
\email{akai@fawstudio.com}
\affiliation{Independent Researcher}

\date{\today}

\begin{abstract}
We prove a no-go theorem connecting time-reversibility, quantum
computation, and the arrow of time within the Petz recovery
framework.  A quantum channel has $\tauasym = 0$ for all input states
if and only if it is unitary (Theorem~\ref{thm:structure}).  It follows
that ideal unitary circuits are simultaneously time-reversible and
computationally universal (Theorem~\ref{thm:ideal}), while any
physical gate with noise breaks time symmetry
(Theorem~\ref{thm:noisy}).  These results yield a trilemma: a physical
system can have at most two of \{time-reversibility, quantum
coherence, environmental interaction\} (Corollary~\ref{cor:trilemma}).
The saturation surface of the Petz recovery bound is
observer-dependent---the ``boundary of the time machine'' depends on
the observer's prior (Theorem~\ref{thm:observer}).  For non-Markovian
dynamics, local time reversal ($d\tauasym/dt < 0$) is possible but
fundamentally bounded by envelope decay and total entropy production.
The past has not disappeared; it is information leaked to the
environment.  $\tauasym$ measures how much has leaked.  Time is the
cost of novelty.
\end{abstract}

\maketitle

%-----------------------------------------------------------------------
\section{Introduction}
\label{sec:intro}
%-----------------------------------------------------------------------

The quantum eraser~\cite{Scully1982,Kim2000} demonstrates that a
future measurement on an idler photon can restore interference in a
signal photon whose which-path information has been recorded.  The
future measurement affects the present interpretation.  But there is
an asymmetry: in an open system, one must wait for the future
measurement to actually occur before the past can be reinterpreted.
The arrow of time appears precisely because the system is open.

The Petz recovery map~\cite{Petz1986,Petz1988} provides the
mathematical structure underlying this observation.  Given a quantum
channel $\chan$ and a reference state $\sigma$, the Petz map
$\Petz_{\sigma,\chan}$ is the unique Bayesian retrodiction
functor~\cite{Parzygnat2023}: it recovers the past from the future.
The recovery infidelity
$\tauasym = 1 - F(\rho, \Petz_{\sigma,\chan}(\chan(\rho)))$ measures
how much the past differs from its retrodicted version---the
operational arrow of time.  When $\tauasym = 0$ for all inputs,
prediction and retrodiction are equivalent, and time has no direction.
When $\tauasym > 0$, some information about the past is gone, and the
arrow of time exists.

In Ref.~\cite{Huang2026a} (hereafter Paper~1), we established the
equivalence chain
$\tauasym = 0 \Leftrightarrow \Sigma = 0 \Leftrightarrow
I(A;E|B) = 0 \Leftrightarrow \text{QMC}$
linking vanishing temporal asymmetry to vanishing entropy production,
vanishing information leakage, and the quantum Markov condition.  We
also proved a composition theorem showing that $\sqrt{\tauasym}$
satisfies a triangle inequality under channel composition.

In this Letter, we use these tools to prove a no-go theorem for ``time
machines.''  The central question is: can a physical system be
simultaneously time-reversible ($\tauasym = 0$) and computationally
useful?  The answer reveals a trilemma connecting time-reversibility,
quantum coherence, and environmental interaction.

The physical picture that emerges is this: the past has not
disappeared.  It is encoded in the environment---in the photons that
scattered off objects, the thermal vibrations of air molecules, the
decoherence records in the surrounding matter.  The parameter
$\tauasym$ measures how much information has leaked.  A ``time
machine,'' in this framework, is not a device that bends spacetime.
It is a device that recovers information from the environment.

%-----------------------------------------------------------------------
\section{Framework}
\label{sec:framework}
%-----------------------------------------------------------------------

We work within the framework of Paper~1~\cite{Huang2026a}.  All
Hilbert spaces are finite-dimensional.

The \emph{temporal asymmetry parameter} for a CPTP map
$\chan: \mathcal{B}(\hilb_A) \to \mathcal{B}(\hilb_B)$, input state
$\rho \in \mathcal{S}(\hilb_A)$, and faithful reference state
$\sigma \in \mathcal{S}(\hilb_A)$ is
\begin{equation}
\tauasym(\rho, \chan, \sigma)
= 1 - F\!\big(\rho,\, \Petz_{\sigma,\chan}(\chan(\rho))\big),
\label{eq:tau_def}
\end{equation}
where $\Petz_{\sigma,\chan}$ is the Petz recovery
map~\cite{Petz1986,Petz1988} and $F$ is the Uhlmann fidelity.

A CPTP map $\chan$ is \emph{time-reversible} if
$\tauasym(\rho, \chan, \sigma) = 0$ for all states $\rho$ and all
faithful $\sigma$.  Equivalently (Paper~1, Eq.~(10)):
$\tauasym = 0$ for all inputs iff $\Sigma = 0$ iff the quantum Markov
condition holds iff perfect Petz recovery is achieved.

A finite set of CPTP maps
$\mathcal{G} = \{\chan_1, \ldots, \chan_k\}$ acting on
$\hilb = (\mathbb{C}^2)^{\otimes n}$ is \emph{computationally
universal} if, for any unitary $U \in U(2^n)$ and any $\epsilon > 0$,
some finite composition
$\chan_{i_m} \circ \cdots \circ \chan_{i_1}$
approximates $U\rho\,U^\dagger$ to within $\epsilon$ in trace norm
for all $\rho$.

A CPTP map is \emph{noiseless} if it is a unitary conjugation
$\chan(\rho) = V\rho\,V^\dagger$; otherwise it is \emph{noisy}.

%-----------------------------------------------------------------------
\section{The No-Go Theorem}
\label{sec:nogo}
%-----------------------------------------------------------------------

\begin{theorem}[Structure of time-reversible channels]
\label{thm:structure}
Let\/ $\chan: \mathcal{B}(\hilb) \to \mathcal{B}(\hilb)$ be a CPTP
map.  Then
$\tauasym(\rho, \chan, \sigma) = 0$ for all states $\rho$ and all
faithful $\sigma$ if and only if\/ $\chan$ is a unitary conjugation:
$\chan(\rho) = U\rho\,U^\dagger$ for some unitary~$U$.
\end{theorem}

\emph{Proof.}---The ``if'' direction is Corollary~1 of
Paper~1~\cite{Huang2026a}: for unitary $\chan_U$, the Petz recovery is
$\Petz_{\sigma,\chan_U}(X) = U^\dagger X\,U$ independent of $\sigma$,
giving $F = 1$ and $\tauasym = 0$.

For ``only if'': $\tauasym = 0$ for all $\rho$ means
$F(\rho, \Petz_{\sigma,\chan}(\chan(\rho))) = 1$ for all $\rho$.  By
the Petz sufficiency theorem~\cite{Petz1988}, this holds for all
$\rho$ if and only if $\chan$ is sufficient for all pairs
$(\rho, \sigma)$, which requires $\chan$ to be reversible.  A
reversible CPTP map on a finite-dimensional algebra is a
$*$-automorphism, which by the Wigner theorem is unitarily
implemented.\hfill$\square$

\begin{theorem}[Ideal computation is time-reversible]
\label{thm:ideal}
Let\/ $\mathcal{G} = \{\chan_1, \ldots, \chan_k\}$ be a gate set
where each $\chan_i$ is a unitary conjugation.  Then:
\begin{enumerate}
\item[\emph{(a)}] Each $\chan_i$ has $\tauasym_i = 0$.
\item[\emph{(b)}] Any composition
  $\chan_{i_m} \circ \cdots \circ \chan_{i_1}$ is unitary, hence has
  $\tauasym_{\mathrm{total}} = 0$.
\item[\emph{(c)}] If the unitaries $\{U_1, \ldots, U_k\}$ form a
  universal gate set (e.g., $\{H, T, \mathrm{CNOT}\}$), then
  $\mathcal{G}$ is computationally universal.
\end{enumerate}
Therefore, ideal quantum computation is simultaneously
time-reversible and computationally universal.
\end{theorem}

\emph{Proof.}---(a) and (b) follow from Corollary~1 of
Paper~1 and the fact that the composition of unitaries is unitary.
(c) is the Solovay--Kitaev theorem~\cite{Kitaev1997,Dawson2006}:
any dense subgroup of $\mathrm{SU}(d)$ provides $\epsilon$-approximate
universality with gate count
$O(\log^c(1/\epsilon))$.\hfill$\square$

\begin{theorem}[Noisy gates break time symmetry]
\label{thm:noisy}
Let $\chan$ be a CPTP map that is not a unitary conjugation.  Then
there exist states $\rho$ and $\sigma$ such that
$\tauasym(\rho, \chan, \sigma) > 0$.
\end{theorem}

\emph{Proof.}---Contrapositive of
Theorem~\ref{thm:structure}.\hfill$\square$

\emph{The naive argument and its failure.}---One might attempt a
different proof: from the saturation theorem (Paper~1, Theorem~4),
$\tauasym = 0$ requires $[\chan(\rho), \chan(\sigma)] = 0$, so a
time-reversible channel maps everything to a commuting subalgebra
and cannot generate quantum coherence.  This argument has a gap.
The saturation theorem characterizes when $F^2 = \exp(-\Delta D)$
is an equality---but $\tauasym = 0$ requires $F = 1$, a stronger
condition.  The correct route is through Theorem~\ref{thm:structure},
which uses the full Petz sufficiency theorem rather than the
saturation condition.

The key point is that Theorem~\ref{thm:ideal} appears to ``destroy''
the no-go theorem: ideal quantum computation has $\tauasym = 0$.
The real no-go emerges from the distinction between mathematical
and physical computation.  Any physical quantum computer operates
with noisy gates.  The physical process implementing gate $U_i$ is
not the ideal unitary conjugation but a noisy CPTP map that
approximates it.  By Theorem~\ref{thm:noisy}, every such gate
contributes $\tauasym > 0$.

The arrow of time in a quantum computer comes from noise, not from
computation.  The unitary content of computation contributes zero
to the time arrow.  Only the irreversible periphery---initialization,
measurement, decoherence---has $\tauasym > 0$.

%-----------------------------------------------------------------------
\section{The Trilemma}
\label{sec:trilemma}
%-----------------------------------------------------------------------

\begin{corollary}[The time--coherence--environment trilemma]
\label{cor:trilemma}
A physical system can achieve at most two of the following three
properties simultaneously:
\begin{enumerate}
\item[\emph{(i)}] \textbf{Time-reversibility:}
  $\tauasym = 0$ for all inputs.
\item[\emph{(ii)}] \textbf{Quantum coherence:}
  non-commuting output states.
\item[\emph{(iii)}] \textbf{Environmental interaction:}
  $\chan$ is not unitary.
\end{enumerate}
Any two imply the negation of the third.
\end{corollary}

\emph{Proof.}---Three cases.

(i) + (ii) $\Rightarrow$ $\lnot$(iii):
Time-reversible quantum coherence (non-commuting outputs with
$\tauasym = 0$ for all inputs) requires, by
Theorem~\ref{thm:structure}, that $\chan$ is unitary---no
environmental interaction.

(i) + (iii) $\Rightarrow$ $\lnot$(ii):
A time-reversible non-unitary channel has
$\tauasym = 0$ for all inputs and is non-unitary.  But
Theorem~\ref{thm:structure} says $\tauasym = 0$ for all inputs
requires unitarity---contradiction unless we weaken to
state-specific $\tauasym = 0$.  A non-unitary channel with
$\tauasym(\rho, \chan, \sigma) = 0$ for specific pairs maps those
states to a commuting subalgebra (Paper~1, Theorem~4), so there is
no quantum coherence between them.

(ii) + (iii) $\Rightarrow$ $\lnot$(i):
Quantum coherence in an open system means $\chan$ is not unitary
and produces non-commuting outputs.  By Theorem~\ref{thm:noisy},
$\tauasym > 0$ for some inputs.\hfill$\square$

Table~\ref{tab:trilemma} summarizes the three scenarios.

\begin{table}[htbp]
\caption{\label{tab:trilemma}The trilemma.  Each row satisfies two of
the three properties and violates the third.  The column ``Physical?''
indicates whether the scenario is realizable in a laboratory.}
\begin{ruledtabular}
\begin{tabular}{lcccl}
\textbf{Scenario} & (i) & (ii) & (iii) & \textbf{Physical?} \\
\colrule
Ideal quantum computer &
  $\checkmark$ & $\checkmark$ & $\times$ &
  No (requires zero noise) \\
Classical open system &
  $\checkmark$ & $\times$ & $\checkmark$ &
  Yes (commuting outputs) \\
Real quantum computer &
  $\times$ & $\checkmark$ & $\checkmark$ &
  Yes (arrow of time) \\
\end{tabular}
\end{ruledtabular}
\end{table}

\emph{Connection to Landauer--Bennett.}---Landauer~\cite{Landauer1961}
showed that erasure costs at least $\kB T \ln 2$ per bit.
Bennett~\cite{Bennett1973} showed that logically reversible
computation can in principle be performed without erasure, at zero
energy cost.  The trilemma is the quantum-information-theoretic
counterpart: unitary computation has $\tauasym = 0$; only the
irreversible periphery (noise, measurement, erasure) has
$\tauasym > 0$.  The logical content of computation is free.  The
physical overhead is what costs.

%-----------------------------------------------------------------------
\section{Observer-Dependent Boundary}
\label{sec:observer}
%-----------------------------------------------------------------------

The saturation surface of the Petz recovery bound (Paper~1,
Theorem~4) divides the space of triples $(\chan, \rho, \sigma)$
into a recoverable region ($\tauasym \approx 0$) and an irreversible
region ($\tauasym \gg 0$).

\begin{theorem}[Observer-dependent time-reversibility]
\label{thm:observer}
For the same channel $\chan$ and the same input state $\rho$, two
observers using different reference states $\sigma_1 \neq \sigma_2$
can disagree on whether the process is time-reversible:
\begin{equation}
\tauasym(\rho, \chan, \sigma_1) = 0
\;\;\not\!\!\!\Rightarrow\;\;
\tauasym(\rho, \chan, \sigma_2) = 0.
\label{eq:observer}
\end{equation}
\end{theorem}

\emph{Proof.}---By example.  Consider the amplitude damping channel
$\chan_\gamma$ with $\gamma = 0.1$ and $\rho = \ket{0}\!\bra{0}$.
Take $\sigma_1 = I/2$.  Then
$\chan(\rho) = \ket{0}\!\bra{0}$ and
$\chan(\sigma_1) = \mathrm{diag}((1+\gamma)/2,\,(1-\gamma)/2)$.
Since $\chan(\rho)$ is rank-1 and $[\chan(\rho), \chan(\sigma_1)] = 0$,
the likelihood ratio is constant on the support, so the saturation
conditions of Paper~1, Theorem~4, are met and $\tauasym$ is small.

Now take $\sigma_2 = \ket{1}\!\bra{1}$.  Then
$\chan(\sigma_2) = \mathrm{diag}(\gamma, 1-\gamma)$, giving a
different likelihood ratio and a different relative entropy gap.  The
recovery fidelity changes, and in general
$\tauasym(\rho, \chan, \sigma_1) \neq
\tauasym(\rho, \chan, \sigma_2)$.\hfill$\square$

The ``boundary of the time machine'' is not absolute.  It depends on
what the observer considers the prior state of the world.  An observer
with a prior $\sigma$ close to $\rho$ (more information) sees a
larger region of time-reversibility.  An observer with less
information sees a smaller region.  This connects to the core result
of Paper~1: the arrow of time is not an absolute property of physics
but emerges from the relationship between the system and the
observer's prior.

In the Bures geometry, $\tauasym = 1 - F$ measures the Bures
distance between $\rho$ and its Petz-recovered version.  The
saturation surface can be visualized as a ``cone'' in Bures space:
processes inside the cone are recoverable, processes outside are
irreversible.  The reference state $\sigma$ determines the orientation
of this cone.

%-----------------------------------------------------------------------
\section{Non-Markovian Time Reversal}
\label{sec:nonmarkov}
%-----------------------------------------------------------------------

For Markovian dynamics (Lindbladian evolution), the temporal asymmetry
satisfies~\cite{Huang2026a}
\begin{equation}
\frac{d\tauasym}{dt}
\sim \frac{\dot{\Sigma}}{2}(1 - \tauasym) \geq 0,
\label{eq:markov_mono}
\end{equation}
where $\dot{\Sigma} \geq 0$ is the instantaneous entropy production
rate.  The arrow of time only advances.

For non-Markovian dynamics, the divisibility condition
fails~\cite{Breuer2016} and $\dot{\Sigma}$ can be temporarily
negative, opening the possibility of
\begin{equation}
\frac{d\tauasym}{dt} < 0
\qquad \text{(local time reversal)}.
\label{eq:nonmarkov}
\end{equation}
Information flows back from the environment to the system, and the
Petz recovery temporarily improves without intervention.

\emph{Jaynes--Cummings model.}---A qubit coupled to a single cavity
mode with Lorentzian spectral density
$J(\omega) = \gamma_0 \lambda^2 /
[2\pi((\omega_0 - \omega)^2 + \lambda^2)]$
provides a concrete example~\cite{Breuer2009}.  The decoherence
function is
\begin{equation}
G(t) = e^{-\lambda t/2}
\!\left[\cosh\!\frac{dt}{2}
+ \frac{\lambda}{d}\sinh\!\frac{dt}{2}\right]\!,
\label{eq:decoherence}
\end{equation}
where $d = \sqrt{\lambda^2 - 2\gamma_0\lambda}$.

In the weak-coupling regime ($\gamma_0/\lambda < 1/2$), $G(t)$ decays
monotonically and $\tauasym(t)$ increases monotonically---standard
Markovian behavior.  In the strong-coupling regime
($\gamma_0/\lambda > 1/2$), $d$ becomes imaginary, $G(t)$ oscillates,
and $\tauasym(t)$ exhibits intervals where
$d\tauasym/dt < 0$.  The fidelity $F(t)$ temporarily increases:
the Petz recovery improves, constituting a concrete
non-Markovianity witness distinct from the BLP trace-distance
witness~\cite{Breuer2009} and the RHP divisibility
witness~\cite{Rivas2010}.

\emph{Design principles.}---Four conditions favor local time reversal:
(1)~a structured environment (spectral peaks, finite bandwidth);
(2)~strong coupling ($\gamma_0/\lambda > 1/2$);
(3)~a finite-dimensional environment (information mirror, not
information sink); and (4)~coherent system--environment interaction
(dephasing-type interactions do not produce backflow).

\emph{Fundamental limits.}---Even with optimal non-Markovian
engineering, three constraints apply.

\emph{Limit~1: Envelope decay.}  The oscillations in $\tauasym(t)$
occur within a decaying envelope:
$\tauasym(t) \sim 1 - |G(t)|^2$ with
$|G(t)| \leq e^{-\lambda t/2}$.  Local time reversal is always
temporary.

\emph{Limit~2: Total entropy production.}  Even when
$d\tauasym/dt < 0$ in the subsystem, the total entropy production
$\Sigma_{\mathrm{total}} \geq 0$ for system plus environment.  The
subsystem's temporary $\tauasym$ decrease is paid for by entropy
increase elsewhere.

\emph{Limit~3: Isolation paradox.}  To achieve $\tauasym \to 0$
permanently requires $\lambda \to 0$: zero dissipation, complete
isolation.  But a completely isolated system undergoes unitary
evolution, and by Theorem~\ref{thm:ideal} can perform quantum
computation internally.  The problem is that one cannot program it
(programming requires interaction: $\tauasym > 0$), cannot read out
the result (measurement requires interaction: $\tauasym > 0$), and
cannot change the computation (the evolution is fixed by the initial
Hamiltonian).  A time machine that cannot be programmed or read out
is not a useful computer.

This closes the loop to the no-go theorem.

%-----------------------------------------------------------------------
\section{The Past as Leaked Information}
\label{sec:past}
%-----------------------------------------------------------------------

The preceding results lead to a unified physical interpretation.

The past has not disappeared.  It is encoded in the
environment---in the photons that scattered off objects, the
thermal vibrations of air molecules, the decoherence records in
surrounding matter.  The parameter $\tauasym$ measures how much
information has leaked from the system to the environment.  A
``time machine,'' in the $\tauasym$ framework, is not a device
that bends spacetime.  It is a device that recovers information
from the environment.

Consider a universe with $\tauasym = 0$ everywhere.  All evolution is
unitary.  All processes are perfectly reversible.  The Petz map
achieves exact recovery everywhere.  Past and future are operationally
indistinguishable.  No information is created or destroyed.  Nothing
happens that was not already determined by initial conditions.

Now consider a universe with $\tauasym > 0$.  Non-unitary processes
exist.  Some information leaks irreversibly to the environment.  The
past is not perfectly recoverable from the future.  Time has a
direction.  Measurement outcomes are genuinely random.  New
information can be created.

The $\tauasym = 0$ universe is a time machine---but it contains no
novelty.  The $\tauasym > 0$ universe has novelty---but it cannot
reverse time.

This tension is quantitative.  The arrow of time $\tauasym$ is bounded
from below by the non-commutativity of the channel outputs (the
``computational content'' of the process) and bounded from above by
the entropy production.  It accumulates sub-additively across
sequential processes via the composition theorem (Paper~1,
Theorem~5).  Time is not an all-or-nothing binary.  It is a continuous
parameter.

\emph{Time is the cost of novelty.}

%-----------------------------------------------------------------------
\section{Discussion}
\label{sec:discussion}
%-----------------------------------------------------------------------

\emph{Trinity of computational bounds.}---The no-go theorem takes its
place alongside two other fundamental constraints on physical
computation.  The Margolus--Levitin bound~\cite{Margolus1998} limits
how \emph{fast} one can compute:
$t_{\min} = \pi\hbar/(2\Delta E)$.  The Landauer
bound~\cite{Landauer1961} limits how \emph{cheaply} one can compute:
$W_{\mathrm{diss}} \geq \kB T \ln 2$ per bit erased.  The Petz--$\tauasym$
bound limits how \emph{reversibly} one can compute:
$\tauasym > 0$ for any noisy gate.  Speed, cost, and reversibility
together define the physical space of realizable computation.

\emph{Connection to ER=EPR.}---Maldacena and
Susskind~\cite{Maldacena2013} conjectured that entanglement is
geometrically dual to wormholes.  Susskind~\cite{Susskind2016} and
Brown et~al.~\cite{Brown2016} connected the growth of the
Einstein--Rosen bridge to the growth of quantum computational
complexity.  In our framework, $\tauasym$ grows as complexity grows:
the wormhole elongates as information leaks irreversibly.  A
``time machine'' ($\tauasym = 0$) corresponds to a static wormhole
with no complexity growth---possible only for an eternal black hole
in perfect thermal equilibrium.  The no-go theorem in this language
becomes: one cannot simultaneously have a growing wormhole (useful
computation) and a static one (time-reversibility).

\emph{Connection to the quantum eraser.}---The quantum eraser is the
physical paradigm from which the entire $\tauasym$ framework
originates.  The eraser shows that a future measurement can restore
past coherence---but only in a closed system, and only for the
specific entangled state.  The $\tauasym$ framework generalizes this:
in a closed system ($\tauasym = 0$), retrodiction is exact and
temporal order is irrelevant.  In an open system ($\tauasym > 0$),
one must wait for the future to actually happen.  The no-go theorem
adds: the price of temporal coexistence ($\tauasym = 0$) is
that the system cannot produce genuinely new information.  In a
$\tauasym = 0$ world, everything is predetermined by initial
conditions.  There are no surprises.

\emph{Outlook.}---The composition theorem (Paper~1, Theorem~5) and the
design principles of Sec.~\ref{sec:nonmarkov} suggest an engineering
program: a ``$\tauasym$-Chrono'' engine that uses Bayesian
composition, saturation scanning, and non-Markovian backflow
optimization to design protocols for quantum memory extension, passive
error correction, and precision tests of time symmetry.

\begin{acknowledgments}
The author thanks M.~M.~Wilde for valuable feedback on the
$\tauasym$ framework, and A.~J.~Parzygnat and F.~Buscemi for
foundational work on retrodiction.
\end{acknowledgments}

%-----------------------------------------------------------------------
% BIBLIOGRAPHY
%-----------------------------------------------------------------------

\begin{thebibliography}{40}

\bibitem{Scully1982}
M.~O. Scully and K.~Dr\"uhl,
``Quantum eraser: A proposed photon correlation experiment concerning
observation and `delayed choice' in quantum mechanics,''
Phys.\ Rev.\ A \textbf{25}, 2208 (1982).

\bibitem{Kim2000}
Y.-H. Kim, R.~Yu, S.~P. Kulik, Y.~Shih, and M.~O. Scully,
``Delayed `choice' quantum eraser,''
Phys.\ Rev.\ Lett.\ \textbf{84}, 1 (2000).

\bibitem{Petz1986}
D.~Petz,
``Sufficient subalgebras and the relative entropy of states of a
von Neumann algebra,''
Commun.\ Math.\ Phys.\ \textbf{105}, 123 (1986).

\bibitem{Petz1988}
D.~Petz,
``Sufficiency of channels over von Neumann algebras,''
Q.\ J.\ Math.\ \textbf{39}, 97 (1988).

\bibitem{Parzygnat2023}
A.~J. Parzygnat and F.~Buscemi,
``Axioms for retrodiction: Achieving time-reversal symmetry with a
prior,''
Quantum \textbf{7}, 1013 (2023).

\bibitem{Huang2026a}
S.-K. Huang,
``The arrow of time from Petz recovery,''
(2026), \url{https://github.com/akaiHuang/petz-recovery-unification}.

\bibitem{Kitaev1997}
A.~Yu.~Kitaev,
``Quantum computations: Algorithms and error correction,''
Russ.\ Math.\ Surv.\ \textbf{52}, 1191 (1997).

\bibitem{Dawson2006}
C.~M. Dawson and M.~A. Nielsen,
``The Solovay--Kitaev algorithm,''
Quantum Inf.\ Comput.\ \textbf{6}, 81 (2006).

\bibitem{Landauer1961}
R.~Landauer,
``Irreversibility and heat generation in the computing process,''
IBM J.\ Res.\ Dev.\ \textbf{5}, 183 (1961).

\bibitem{Bennett1973}
C.~H. Bennett,
``Logical reversibility of computation,''
IBM J.\ Res.\ Dev.\ \textbf{17}, 525 (1973).

\bibitem{ReebWolf2014}
D.~Reeb and M.~M. Wolf,
``An improved Landauer principle with finite-size corrections,''
New J.\ Phys.\ \textbf{16}, 103011 (2014).

\bibitem{Margolus1998}
N.~Margolus and L.~B. Levitin,
``The maximum speed of dynamical evolution,''
Physica D \textbf{120}, 188 (1998).

\bibitem{Lloyd2000}
S.~Lloyd,
``Ultimate physical limits to computation,''
Nature \textbf{406}, 1047 (2000).

\bibitem{Junge2018}
M.~Junge, R.~Renner, D.~Sutter, M.~M. Wilde, and A.~Winter,
``Universal recovery maps and approximate sufficiency of quantum
relative entropy,''
Ann.\ Henri Poincar\'e \textbf{19}, 2955 (2018).

\bibitem{Fawzi2015}
O.~Fawzi and R.~Renner,
``Quantum conditional mutual information and approximate Markov
chains,''
Commun.\ Math.\ Phys.\ \textbf{340}, 575 (2015).

\bibitem{Kolchinsky2017}
A.~Kolchinsky and D.~H. Wolpert,
``Dependence of dissipation on the initial distribution over states,''
J.\ Stat.\ Mech.\ \textbf{2017}, 083202 (2017).

\bibitem{Breuer2009}
H.-P. Breuer, E.-M. Laine, and J.~Piilo,
``Measure for the degree of non-Markovian behavior of quantum
processes in open systems,''
Phys.\ Rev.\ Lett.\ \textbf{103}, 210401 (2009).

\bibitem{Rivas2010}
A.~Rivas, S.~F. Huelga, and M.~B. Plenio,
``Entanglement and non-Markovianity of quantum evolutions,''
Phys.\ Rev.\ Lett.\ \textbf{105}, 050403 (2010).

\bibitem{Breuer2016}
H.-P. Breuer, E.-M. Laine, J.~Piilo, and B.~Vacchini,
``Colloquium: Non-Markovian dynamics in open quantum systems,''
Rev.\ Mod.\ Phys.\ \textbf{88}, 021002 (2016).

\bibitem{Crooks1999}
G.~E. Crooks,
``Entropy production fluctuation theorem and the nonequilibrium
work relation for free energy differences,''
Phys.\ Rev.\ E \textbf{60}, 2721 (1999).

\bibitem{Kwon2019}
H.~Kwon and M.~S. Kim,
``Fluctuation theorems for a quantum channel,''
Phys.\ Rev.\ X \textbf{9}, 031029 (2019).

\bibitem{Maldacena2013}
J.~Maldacena and L.~Susskind,
``Cool horizons for entangled black holes,''
Fortschr.\ Phys.\ \textbf{61}, 781 (2013).

\bibitem{Susskind2016}
L.~Susskind,
``Computational complexity and black hole horizons,''
Fortschr.\ Phys.\ \textbf{64}, 24 (2016).

\bibitem{Brown2016}
A.~R. Brown, D.~A. Roberts, L.~Susskind, B.~Swingle, and Y.~Zhao,
``Complexity, action, and black holes,''
Phys.\ Rev.\ D \textbf{93}, 086006 (2016).

\bibitem{Zurek2003}
W.~H. Zurek,
``Decoherence, einselection, and the quantum origins of the
classical,''
Rev.\ Mod.\ Phys.\ \textbf{75}, 715 (2003).

\bibitem{Zeh1970}
H.~D. Zeh,
``On the interpretation of measurement in quantum theory,''
Found.\ Phys.\ \textbf{1}, 69 (1970).

\bibitem{Landi2021}
G.~T. Landi and M.~Paternostro,
``Irreversible entropy production: From classical to quantum,''
Rev.\ Mod.\ Phys.\ \textbf{93}, 035008 (2021).

\bibitem{Buscemi2024}
G.~Bai, F.~Buscemi, and V.~Scarani,
``Fully quantum stochastic entropy production,''
arXiv:2412.12489 (2024).

\bibitem{Barnum2002}
H.~Barnum and E.~Knill,
``Reversing quantum dynamics with near-optimal quantum and classical
fidelity,''
J.\ Math.\ Phys.\ \textbf{43}, 2097 (2002).

\bibitem{Fullwood2023}
A.~J. Parzygnat and J.~Fullwood,
``From time-reversal symmetry to quantum Bayes' rules,''
PRX Quantum \textbf{4}, 020334 (2023).

\bibitem{Png2025}
W.-H. Png and V.~Scarani,
``Petz recovery maps of single-qubit decoherence channels in an
ion trap quantum processor,''
Phys.\ Rev.\ A \textbf{112}, 022613 (2025).

\bibitem{Singh2025}
G.~Singh, R.~S. Sahani, V.~Jagadish, L.~Lautenbacher,
N.~K. Bernardes, and K.~Dorai,
``Realizing the Petz recovery map on an NMR quantum processor,''
arXiv:2508.08998 (2025).

\bibitem{ABL1964}
Y.~Aharonov, P.~G. Bergmann, and J.~L. Lebowitz,
``Time symmetry in the quantum process of measurement,''
Phys.\ Rev.\ \textbf{134}, B1410 (1964).

\end{thebibliography}

\end{document}
